\documentclass[letterpaper,10pt]{article}
\usepackage{latexsym}
\usepackage[empty]{fullpage}
\usepackage{titlesec}
\usepackage{marvosym}
\usepackage[usenames,dvipsnames]{color}
\usepackage{verbatim}
\usepackage{enumitem}
\usepackage[hidelinks]{hyperref}
\usepackage{fancyhdr}
\usepackage[english]{babel}
\usepackage{tabularx}

\pagestyle{fancy}
\fancyhf{} % clear all header and footer fields
\fancyfoot{}
\renewcommand{\headrulewidth}{0pt}
\renewcommand{\footrulewidth}{0pt}

% Adjust margins
\addtolength{\oddsidemargin}{-0.5in}
\addtolength{\evensidemargin}{-0.5in}
\addtolength{\textwidth}{1in}
\addtolength{\topmargin}{-.85in}
\addtolength{\textheight}{1.75in}

\urlstyle{same}

\raggedbottom
\raggedright
\setlength{\tabcolsep}{0in}

% Sections formatting
\titleformat{\section}{
  \vspace{-7pt}\scshape\raggedright\large
}{}{0em}{}[\color{black}\titlerule \vspace{-5pt}]

% Ensure that generated pdf is machine readable/ATS parsable
\ifdefined\pdfgentounicode
  \input{glyphtounicode}
  \pdfgentounicode=1
\fi

%-------------------------
% Custom commands
\newcommand{\resumeItem}[1]{
  \item\small{
    {#1 \vspace{-2pt}}
  }
}

\newcommand{\resumeSubheading}[4]{
  \vspace{-1pt}\item
    \begin{tabular*}{0.97\textwidth}[t]{l@{\extracolsep{\fill}}r}
      \textbf{#1} & #2 \\
      \textit{\small#3} & \textit{\small #4} \\
    \end{tabular*}\vspace{-5pt}
}

\newcommand{\resumeSubSubheading}[2]{
    \item
    \begin{tabular*}{0.97\textwidth}{l@{\extracolsep{\fill}}r}
      \textit{\small#1} & \textit{\small #2} \\
    \end{tabular*}\vspace{-5pt}
}

\newcommand{\resumeProjectHeading}[2]{
    \item
    \begin{tabular*}{0.97\textwidth}{l@{\extracolsep{\fill}}r}
      \small#1 & #2 \\
    \end{tabular*}\vspace{-5pt}
}

\newcommand{\resumeSubItem}[1]{\resumeItem{#1}\vspace{-3pt}}

\renewcommand\labelitemii{$\vcenter{\hbox{\tiny$\bullet$}}$}

\newcommand{\resumeSubHeadingListStart}{\begin{itemize}[leftmargin=0.15in, label={}]}
\newcommand{\resumeSubHeadingListEnd}{\end{itemize}}
\newcommand{\resumeItemListStart}{\begin{itemize}}
\newcommand{\resumeItemListEnd}{\end{itemize}\vspace{-3pt}}

%-------------------------------------------
%%%%%%  RESUME STARTS HERE  %%%%%%%%%%%%%%%%%%%%%%%%%%%%


\begin{document}

%----------HEADING----------
\begin{center}
    \textbf{\Huge \scshape Abhijit Ghosh} \\ \vspace{3pt}
    \small Mysuru, Karnataka, India - 570016 $|$ (+91) 82938 16687 $|$ \href{mailto:abhijitghoshprem@gmail.com}{\underline{abhijitghoshprem@gmail.com}} $|$
    \href{https://linkedin.com/in/abhijitmechie/}{\underline{linkedin.com/in/abhijitmechie/}} $|$
    \href{https://abhijittorq.vercel.app/}{\underline{Portfolio}}
\end{center}
\vspace{-12pt}

%-----------OBJECTIVE-----------
\section{Objective}
\small{Geometrical Assurance \& Mechanical Design Engineer with expertise in 3D tolerance stack-up analysis (CETOL), CAD modeling (Siemens NX), and digital mockup validation (Teamcenter Vis Mockup). Applying ASME Y14.5 GD\&T and 3-2-1 datum strategies across engine subsystems to eliminate packaging clashes and optimize assembly yield.}
\vspace{-5pt}

%-----------EXPERIENCE-----------
\section{Professional Experience}
  \resumeSubHeadingListStart
    \resumeSubheading
      {Senior Systems Engineer (Stellantis Program)}{October 2026 -- Present}
      {Infosys Limited}{Mysuru, India}
      \resumeItemListStart
        \resumeItem{Increased powertrain subsystem assembly yield from 54.34\% to 99.90\% per million units by executing 3D variation simulation in CETOL and reallocating critical GD\&T tolerances.}
        \resumeItem{Built automated project tracking workflows using Microsoft Power Automate and Lists, cutting administrative reporting overhead and enabling real-time engineering change tracking with Stellantis stakeholders.}
        \resumeItem{Promoted to Senior Systems Engineer in recognition of consistently high performance and engineering impact on the Stellantis program.}
      \resumeItemListEnd

    \resumeSubheading
      {Systems Engineer (Stellantis Program)}{May 2025 -- September 2026}
      {Infosys Limited}{Mysuru, India}
      \resumeItemListStart
        \resumeItem{Conducted 3D tolerance stack-up analysis in CETOL and Siemens NX across engine cooling circuits and air intake packaging, mitigating high-risk assembly variation before tooling freeze.}
        \resumeItem{Performed dynamic clash and gap/flush analyses in Teamcenter Vis Mockup, resolving clearance violations across critical powertrain packaging interfaces.}
        \resumeItem{Applied ASME Y14.5 GD\&T standards and 3-2-1 datum strategies in cross-functional design reviews to ensure dimensional repeatability and prevent build defects.}
      \resumeItemListEnd

    \resumeSubheading
      {Systems Engineer Trainee}{October 2024 -- April 2025}
      {Infosys Limited}{Mysuru, India}
      \resumeItemListStart
        \resumeItem{Evaluated component fit and alignment across early-stage assemblies using ASME Y14.5 GD\&T and 1D tolerance stack-up calculations.}
        \resumeItem{Developed 3D models and manufacturing detail drawings in Siemens NX and CATIA V5 to support design modifications and manufacturing readiness.}
        \resumeItem{Performed structural FEA in Ansys to verify stress distributions and validate design assumptions.}
      \resumeItemListEnd

    \resumeSubheading
      {Subject Matter Expert - Mechanical Engineering}{July 2023 -- September 2024}
      {Vaidik Eduservices Pvt. Ltd}{Remote, India}
      \resumeItemListStart
        \resumeItem{Delivered technical problem-solving and rigorous engineering solutions across core mechanical engineering domains (thermodynamics, mechanics of materials, and machine design).}
      \resumeItemListEnd

    \resumeSubheading
      {Intern}{March 2023 -- June 2023}
      {IIT Jodhpur - Dept of Metallurgical and Materials Engineering}{Jodhpur, India}
      \resumeItemListStart
        \resumeItem{Investigated novel in-situ volume contractible metal halide negative electrodes for high-performance lithium-ion batteries, executing material characterization and experimental data analysis.}
      \resumeItemListEnd
  \resumeSubHeadingListEnd

%-----------SKILLS-----------
\section{Skills}
 \begin{itemize}[leftmargin=0.15in, label={}]
    \small{\item{
     \textbf{Core CAD / Simulation}{: Siemens NX, CETOL (3D VSA / Tolerance Stack-Up), Teamcenter Vis Mockup, CATIA V5, Ansys, Ansys Workbench, Product \& Manufacturing Information (PMI), Open CAD} \\
     \textbf{Standards \& Methods}{: ASME Y14.5 (GD\&T), 3-2-1 Datum Strategy, Gap \& Flush Analysis, 1D/3D Tolerance Analysis, DFMEA} \\
     \textbf{PLM \& Data}{: Siemens Teamcenter, NX PLM, SQL} \\
     \textbf{Automation \& Tools}{: Microsoft Power Automate, Microsoft Lists, Microsoft Power Platform, Java} \\
     \textbf{Languages}{: English, German (Elementary Proficiency)}
    }}
 \end{itemize}
 \vspace{-15pt}

%-----------EDUCATION-----------
\section{Education}
  \resumeSubHeadingListStart
    \resumeSubheading
      {Techno Main Salt Lake}{Kolkata, India}
      {Bachelor of Technology in Mechanical Engineering -- \textbf{DGPA: 8.55}}{July 2022}
    \resumeSubheading
      {Santiniketan Institute of Polytechnic}{Bolpur, India}
      {Diploma in Mechanical Engineering -- \textbf{DGPA: 7.8}}{June 2019}
  \resumeSubHeadingListEnd

%-----------CERTIFICATIONS-----------
\section{Certifications}
 \begin{itemize}[leftmargin=0.15in, label={}]
    \item \small{Certified Lean Six Sigma AI Yellow Belt -- Sparen \& Gewinn Consulting (Aug 2026)}
    \vspace{-6pt}
    \item \small{DFMEA-MSR: Design FMEA \& Monitoring System Response (Dec 2025)}
    \vspace{-6pt}
    \item \small{Geometric Dimensioning \& Tolerancing (GD\&T) -- Infosys (Mar 2025)}
    \vspace{-6pt}
    \item \small{NX PLM-1 Engineering Author -- Infosys (May 2026)}
    \vspace{-6pt}
    \item \small{Teamcenter 14.1x Visualization Mockup -- Tata Technologies (Jan 2026)}
    \vspace{-6pt}
    \item \small{Product \& Manufacturing Information (PMI) -- Tata Technologies (Dec 2025)}
    \vspace{-6pt}
    \item \small{Fundamentals of Finite Element Analysis -- Tata Technologies (Jan 2026)}
    \vspace{-6pt}
    \item \small{SOLIDWORKS Certification -- Academy of Skill Dev (Mar 2021)}
    \vspace{-6pt}
    \item \small{Design for Additive Manufacturing -- Infosys (Jul 2025)}
 \end{itemize}

%-----------AWARDS-----------
\section{Awards}
 \begin{itemize}[leftmargin=0.15in, label={}]
    \item \small{\textbf{Infosys Insta Award (Feb 2026):} Recognized by leadership for outstanding contributions to project efficiency through custom workflow automation initiatives.}
    \vspace{-4pt}
    \item \small{\textbf{High Performer} in Infosys Foundation Program}
 \end{itemize}

\end{document}
